%% Font size %%
\documentclass[11pt]{article}

%% Load the custom package
\usepackage{Mathdoc}

%% Numéro de séquence %% Titre de la séquence %%
\renewcommand{\centerhead}{S01 - Fractions - Comparaison à l'unité}

%% Spacing commands %%
\renewcommand{\baselinestretch}{1} \setlength{\parindent}{0pt}

\begin{document}

\begin{exercice}[1][Exercice corrigé]
Pour comparer une fraction à $1$, on compare son numérateur et son dénominateur~:
\begin{itemize}[itemsep=.3em]
\item si le numérateur est \textbf{plus petit} que le dénominateur, la fraction est \textbf{plus petite que $1$} ;
\item si le numérateur est \textbf{plus grand} que le dénominateur, la fraction est \textbf{plus grande que $1$} ;
\item si le numérateur est \textbf{égal} au dénominateur, la fraction est \textbf{égale à $1$}.
\end{itemize}

\textbf{Exemples :}
\begin{itemize}[itemsep=.3em]
\item $\dfrac{5}{6}$ : le numérateur ($5$) est plus petit que le dénominateur ($6$), donc $\dfrac{5}{6} < 1$.
\item $\dfrac{8}{3}$ : le numérateur ($8$) est plus grand que le dénominateur ($3$), donc $\dfrac{8}{3} > 1$.
\end{itemize}

\vspace{.5em}
\textbf{À toi :} pour chaque fraction, compare le numérateur et le dénominateur, puis complète avec $<$, $>$ ou $=$.
\begin{multicols}{2}
\begin{enumerate}[itemsep=1em]
\item $\dfrac{3}{4} \ \ldots \ 1$ \dotfill
\item $\dfrac{7}{2} \ \ldots \ 1$ \dotfill
\item $\dfrac{9}{9} \ \ldots \ 1$ \dotfill
\item $\dfrac{4}{11} \ \ldots \ 1$ \dotfill
\end{enumerate}
\end{multicols}
\end{exercice}

\begin{exercice}[1][Comparaison à l'unité]
Complète avec $<$, $>$ ou $=$.
\begin{multicols}{2}
\begin{enumerate}[itemsep=1em]
\item $\dfrac{2}{5} \ \ldots \ 1$
\item $\dfrac{11}{4} \ \ldots \ 1$
\item $\dfrac{6}{6} \ \ldots \ 1$
\item $\dfrac{9}{13} \ \ldots \ 1$
\item $\dfrac{17}{5} \ \ldots \ 1$
\item $\dfrac{8}{15} \ \ldots \ 1$
\end{enumerate}
\end{multicols}
\end{exercice}

\begin{exercice}[3][Avec de grands dénominateurs]
Complète avec $<$, $>$ ou $=$ (attention, les dénominateurs sont supérieurs à $100$).
\begin{multicols}{2}
\begin{enumerate}[itemsep=1em]
\item $\dfrac{87}{124} \ \ldots \ 1$
\item $\dfrac{205}{198} \ \ldots \ 1$
\item $\dfrac{150}{150} \ \ldots \ 1$
\item $\dfrac{99}{101} \ \ldots \ 1$
\item $\dfrac{310}{207} \ \ldots \ 1$
\item $\dfrac{128}{512} \ \ldots \ 1$
\end{enumerate}
\end{multicols}
\end{exercice}

\end{document}
